# Automated Analysis Pipeline Exploration - Cervical Cancer Risk Factors

## Purpose of the Notebook

This notebook runs and extends the instructor's automated analysis pipeline on my final
project dataset, the UCI Cervical Cancer (Risk Factors) dataset. The pipeline loads the
data, checks and prepares it, explores it, runs inferential statistics, and trains and
evaluates supervised machine learning models. It uses conditional logic to adapt to the
file type, missing values, and user-defined settings.

The prediction task is to classify whether a patient's cervical biopsy is positive
(`Biopsy` = 1) from age, sexual and reproductive history, contraceptive use, smoking, and
STD history. It supports risk-stratification ideas only and does not replace diagnostic
testing.

The notebook also explains the pipeline stages and conditional logic, documents the
modifications that were made, and interprets the results of the modified pipeline.

### Modifications

1. **Data loading:** loads from a local file, then GitHub, then (Colab only) an upload prompt, and converts `?` to missing. The notebook runs with no manual upload, and the raw file marks missing answers with `?`.
2. **Automatic cleaning:** fixes impossible values, drops columns that are too empty or constant, and drops duplicate rows. 
3. **Class weights and cutoff:** balanced class weights, and a probability cutoff chosen to reach a target recall.
4. **Random forest and PR AUC:** adds a random forest model and a PR AUC metric. 

Also, the plot variables were changed to fit this dataset and the unseeded random sample was fixed so the output is the same on every run.

## Repository Contents

| File | Description |
|---|---|
| `Week3_AutomatedAnalysisPipelineNotebook.ipynb` | The completed Google Colab notebook (already run, with outputs) |
| `risk_factors_cervical_cancer.csv` | The project dataset (original UCI format is CSV) |
| `README.md` | This file |

## Instructions for Running the Analysis

Open the notebook with the link below, then choose **Runtime > Restart session and run all**.

**Colab link:** https://colab.research.google.com/drive/1-Xt2PyUdwyESYsLgimoGmcETXUY7fjW- 

No upload is needed. The first code cell loads the data in this order:
1. `risk_factors_cervical_cancer.csv` from the current folder, if present.
2. The same file from this GitHub repository (set by `DATA_URL` in the first code cell, which points to `katerina-smith/Programming_for_HDS_Week3b`).
3. If both fail and the notebook is running in Colab, a prompt to upload the file.

The notebook uses fixed random seeds (`random_state=42`), so results are the same on every run.

## Required Software and Libraries

Python 3 with: pandas, numpy, scipy, scikit-learn, matplotlib, seaborn, statsmodels, and phik
(see `requirements.txt`). All except `phik` are pre-installed in Google Colab; the notebook
installs `phik` automatically in the Phi-K correlation cell.

## Dataset

**Cervical Cancer (Risk Factors)**, UCI Machine Learning Repository:
<https://archive.ics.uci.edu/dataset/383/cervical+cancer+risk+factors>
(Hospital Universitario de Caracas, Venezuela, 2017; license CC BY 4.0).
858 patients and 36 columns in the original file. After the automatic cleaning in the
notebook: 835 patients and 32 columns (28 predictors after setting aside the target and
three other screening-test columns).

## Assumptions and Limitations

- **Target and leakage:** the target is `Biopsy`. The columns `Hinselmann`, `Schiller`, and
  `Citology` are other screening-test results that closely track the biopsy result, so they
  are excluded from the predictors.
- **Few positive cases:** about 6 percent of patients are biopsy positive. The test set has only 16
  positive patients, so performance numbers are uncertain and would change with a different
  split. The notebook uses a single train/test split.
- **Missing data:** missing answers are imputed with training-set medians. They are not
  random (patients could skip questions), so imputation cannot recover skipped answers.
- **Statistical tests:** several tests were run without multiple-comparison adjustment, so
  p-values are descriptive.
- **Weak predictive signal:** models perform only modestly. A cutoff chosen to catch 80 percent of
  positive biopsies flags most patients, so these models are not ready for clinical use.
- **Single hospital, self-reported, observational data:** results may not generalize, and
  relationships are associations, not causes.
  
  ## AI Use
  
 Anthropic. Claude. Used to improve documentation, give examples of optimal 
analysis pipeline, troubleshoot/optimize code, and assist with model interpretation. 
Also used to generate examples of README documentation in GitHub.
